% ============================================================
%  Chapter 4 — Servers and Data Centres
%  The Secret Life of the Internet
% ============================================================

\chapter{Servers and Data Centres}

\chapteropening{%
  Every website, video, and game lives somewhere.
  But where?  Come and visit the giant buildings that hold it all!
}
\chapterbadge{Welcome to Server City}

% --- Opening story ---
\section*{The Giant Warehouse of the Internet}

Dot had zoomed through the Wi-Fi signal, past the router,
and along a cable into the bigger internet.

Up ahead was a colossal building, almost the size of several football fields,
humming with the sound of thousands of fans.
The lights blinked in long rows, like a city skyline made of computers.

\character{Dot}{This is a \textbf{data centre} — one of the busiest places on the internet!
  Somewhere inside, a friendly helper called Sunny Server is waiting for us.}

\sectionrule

% --- What is a server? ---
\section*{What Is a Server?}

A \term{server} is a powerful computer that stores information and sends it to people who ask for it.

\begin{funfact}
Google's data centres process over 8.5 billion searches every single day. That's about 99,000 searches every second!
\end{funfact}

When you open a website, your browser sends a request to a server.
The server finds all the pieces of that website — the words, the pictures, the buttons, the videos —
and sends them back to your browser.

\begin{picturethis}
  Imagine a huge library.
  The librarian (the server) knows where every book (piece of information) is kept.
  When you ask for a book, the librarian fetches it and brings it to you.
  The server is the librarian.
  The books are the information.
  You are the person asking.
\end{picturethis}

Inside the data centre, Dot found a row of tall metal shelves.
Each shelf held many flat computers stacked on top of each other.
One of them glowed a warm golden colour.

\character{Sunny Server}{Hello!  I store billions of web pages, photos, and videos.
  When someone asks for something I have,
  I pack it up and send it straight to them — as fast as I can!}

\sectionrule

% --- What is a Data Centre? ---
\section*{What Is a Data Centre?}

A \term{data centre} is a large building that holds hundreds or even thousands of servers.

Data centres are carefully designed to keep all those servers running smoothly:

\begin{itemize}
  \item \textbf{Cooling systems} — Servers get very hot when they work hard, so powerful fans and air conditioners keep the temperature just right.
  \item \textbf{Reliable power} — data centres have backup generators so that if the electricity goes out, the servers keep running.
  \item \textbf{Security} — Only authorised people can get inside. Many have cameras, locks, and security guards.
  \item \textbf{Lots of cables} — Miles and miles of cables connect the servers to each other and to the internet.
\end{itemize}

\begin{picturethis}
  Picture an apartment building — but instead of flats with people,
  every flat is a server computer.
  The building manager (the data centre team) makes sure the heating works,
  the electricity stays on, and the right people have access.
  That building is the data centre.
\end{picturethis}

\sectionrule

% --- Request and response ---
\section*{Fetch and Send}

\begin{quickmap}
  \textbf{What happens when you load a web page:}
  \begin{enumerate}[label=\textbf{\arabic*.}]
    \item Your browser sends a \term{request}: ``Please send me this page!''
    \item The request travels through Wi-Fi, the router, your ISP, and the internet to the server.
    \item Sunny Server finds all the parts of that page:
          the text, images, and layout instructions.
    \item The server packs them up and sends a \term{response} back.
    \item The response travels back the same way to your device.
    \item Your browser assembles all the parts and shows you the finished page.
  \end{enumerate}
\end{quickmap}

\begin{diagram}[Request and response travelling between you and a distant server.]
  \resizebox{\linewidth}{!}{%
  \begin{tikzpicture}[>=Stealth, node distance=1.15cm]
    \node[draw=InternetBlue, rounded corners, fill=SkyBlue!20, minimum width=2.2cm, minimum height=0.9cm] (d) {Your Device};
    \node[draw=WarmOrange!80!black, rounded corners, fill=WarmOrange!20, minimum width=1.8cm, minimum height=0.9cm, right=of d] (r) {Router};
    \node[draw=LeafGreen!80!black, rounded corners, fill=LeafGreen!20, minimum width=1.4cm, minimum height=0.9cm, right=of r] (i) {ISP};
    \node[draw=SoftPurple!80!black, rounded corners, fill=SoftPurple!20, minimum width=1.9cm, minimum height=0.9cm, right=of i] (n) {Internet};
    \node[draw=CoralRed!80!black, rounded corners, fill=CoralRed!16, minimum width=2.6cm, minimum height=0.9cm, right=of n] (s) {Data Centre / Server};

    \draw[<->,thick,InternetBlue] (d) -- (r);
    \draw[<->,thick,InternetBlue] (r) -- (i);
    \draw[<->,thick,InternetBlue] (i) -- (n);
    \draw[<->,thick,InternetBlue] (n) -- (s);

    \node[font=\small\bfseries\color{InternetBlue}] at ($(d)!0.5!(s)+(0,0.8)$) {request $\rightarrow$};
    \node[font=\small\bfseries\color{WarmOrange}] at ($(d)!0.5!(s)+(0,-0.8)$) {$\leftarrow$ response};
  \end{tikzpicture}%
  }
\end{diagram}

All of this happens in less than a second!

\sectionrule

% --- Big companies and their data centres ---
\section*{Who Runs These Giant Buildings?}

Many big companies around the world run data centres.
Large video and search websites have data centres on multiple continents,
so that no matter where you are in the world,
there is a server nearby that can answer your request quickly.

The closer the server is to you, the faster the response.
That is why some videos load almost immediately even from the other side of the world —
there is likely a copy stored on a server much closer to your home.

\sectionrule

\begin{newwords}
  \begin{description}[labelwidth=3.5cm, leftmargin=3.7cm]
    \item[\term{Server}]       A powerful computer that stores information and sends it to people who ask.
    \item[\term{Data centre}]  A large building full of servers, kept cool, powered, and secure.
    \item[\term{Request}]      A message from your browser asking a server for information.
    \item[\term{Response}]     The answer from the server, containing the web page or data you asked for.
    \item[\term{Cache}]        A nearby copy of data stored to help deliver it faster.
  \end{description}
\end{newwords}

\sectionrule

\begin{tryit}
  \textbf{Help Sunny Server pack the right pieces!}

  A visitor has asked for a web page about penguins.
  Circle all the things Sunny Server should pack and send:

  \begin{center}
    \begin{tabularx}{\linewidth}{Y Y Y}
      A photo of a penguin & A penguin fact paragraph & A recipe for pasta \\
      A video of penguins & A headline: ``Meet the Penguin!'' & A map of Antarctica \\
      A button: ``Learn More'' & A picture of a car & A penguin sounds clip \\
    \end{tabularx}
  \end{center}

  \emph{Everything related to penguins gets packed — the rest stays on the shelf!}
\end{tryit}

\sectionrule

\begin{bigidea}
  \textbf{Servers} are powerful computers that store and send information.
  They live in giant buildings called \textbf{data centres},
  which keep them cool, powered, and ready to answer millions of requests every second.
\end{bigidea}

\character{Sunny Server}{Next stop is even more amazing — the cables under the ocean!
  Some of my replies travel thousands of kilometres under the sea to reach you.
  Captain Cable will show you how!}
